\section*{Bab \ref{ch:b1:series} --- Deret Numerik}

\begin{solution}{exo:b1:series:1}
$\dfrac{1}{n(n+1)} = \dfrac1n - \dfrac{1}{n+1}$: jadi teleskopis, dengan $S_N =
1 - \frac{1}{N+1} \to 1$. Konvergen, dengan jumlah $1$.

$\ln\bigl(1 - \frac{1}{n^2}\bigr) = \ln\frac{(n-1)(n+1)}{n^2} =
\ln\frac{n-1}{n} - \ln\frac{n}{n+1}$: jadi teleskopis lagi, dengan $S_N =
\ln\frac12 - \ln\frac{N}{N+1} \to -\ln 2$. Konvergen, dengan jumlah $-\ln 2$.

$\dfrac{3^n + 4^n}{5^n} = \bigl(\frac35\bigr)^n +
\bigl(\frac45\bigr)^n$: yaitu dua \omterm{ex:b1:series:geometric}{deret geometri} yang konvergen, dengan jumlah
$\frac{1}{1 - 3/5} + \frac{1}{1 - 4/5} = \frac52 + 5 = \frac{15}{2}$.
\end{solution}

\begin{solution}{exo:b1:series:2}
Lewat \omterm{thm:b1:series:ratio}{uji nisbahnya} di sepanjangnya.

$\frac{u_{n+1}}{u_n} = \frac{(n+1)^2}{2n^2} \to \frac12 < 1$:
jadi konvergen.

$\frac{u_{n+1}}{u_n} = \frac{(n+1)!\,n^n}{n!\,(n+1)^{n+1}} =
\bigl(\frac{n}{n+1}\bigr)^n = \bigl(1 + \frac1n\bigr)^{-n} \to
\frac1\eu < 1$: jadi konvergen.

Dengan faktor $2^n$: nisbahnya $\to \frac2\eu < 1$: jadi konvergen.

Dengan $3^n$: nisbahnya $\to \frac3\eu > 1$: jadi divergen (karena sukunya menuju
$+\infty$).
\end{solution}

\begin{solution}{exo:b1:series:3}
Semuanya bersuku taknegatif; jadi pakailah kesetaraannya
(\cref{thm:b1:series:positive}).

$\sin\frac{1}{n^2} \sim \frac{1}{n^2}$: jadi konvergen (menurut Riemann dengan $\alpha
= 2$).

$1 - \cos\frac1n \sim \frac{1}{2n^2}$: jadi konvergen.

$\frac{1}{\sqrt{n(n+1)}} \sim \frac1n$: jadi divergen.

$\frac{\ln n}{n^2} = \frac{1}{n^{3/2}}\cdot\frac{\ln n}{n^{1/2}}$
dan $\frac{\ln n}{\sqrt n} \to 0$
(\cref{prop:b1:functions:powerrules}): sehingga $\frac{\ln n}{n^2} \leq
\frac{1}{n^{3/2}}$ untuk $n$ yang besar: jadi konvergen.
\end{solution}

\begin{solution}{exo:b1:series:4}
Untuk $n \geq 2$: $\frac{1}{n^2} \leq \frac{1}{n(n-1)} =
\frac{1}{n-1} - \frac1n$. Sehingga
\[
\sum_{n=1}^{N} \frac{1}{n^2} \leq 1 + \sum_{n=2}^{N}
\Bigl(\frac{1}{n-1} - \frac1n\Bigr) = 1 + 1 - \frac1N < 2 :
\]
jadi jumlah parsialnya naik dan terbatas oleh $2$: sehingga konvergen
(\cref{thm:b1:series:positive}), dengan jumlah $\leq 2$. (Adapun nilai persisnya
$\frac{\pi^2}{6}$ merupakan perayaan tahun kedua.)
\end{solution}

\begin{solution}{exo:b1:series:5}
$\sum \frac{(-1)^n}{\sqrt n}$: yaitu berselang-seling dengan $\frac{1}{\sqrt n}
\downarrow 0$: jadi konvergen (\cref{thm:b1:series:alternating}); tetapi tak
mutlak (karena $\alpha = \frac12 \leq 1$).

$\sum \frac{(-1)^n}{n + (-1)^n}$: barisan $\frac{1}{n +
(-1)^n}$ \emph{tak} menurun (karena $\frac{1}{n+1}$ lalu
$\frac{1}{n}$ berselang-seling dengan buruk), sehingga ujinya tak berlaku secara langsung.
Jadi ekspansikanlah:
\[
\frac{(-1)^n}{n + (-1)^n}
= \frac{(-1)^n}{n}\cdot\frac{1}{1 + \frac{(-1)^n}{n}}
= \frac{(-1)^n}{n} - \frac{1}{n^2} +
O\Bigl(\frac{1}{n^3}\Bigr):
\]
maka \omterm{def:b1:series:def}{deret} pertamanya konvergen (menurut berselang-selingnya), $\sum \frac{1}{n^2}$
konvergen, dan $O(n^{-3})$-nya konvergen mutlak: sehingga jumlah ketiga
\omterm{def:b1:series:def}{deret} yang konvergen itu konvergen.

Untuk $\sin\bigl(\pi\sqrt{n^2+1}\bigr)$: tulislah $\sqrt{n^2 + 1} = n +
\frac{1}{2n} + \varepsilon_n$ dengan $\varepsilon_n = O(n^{-3})$;
lalu, menurut keperiodikan-$\pi$ $\sin$ sampai tandanya,
\[
\sin\bigl(\pi\sqrt{n^2+1}\bigr)
= (-1)^n \sin\Bigl(\frac{\pi}{2n} + \pi\varepsilon_n\Bigr) .
\]
Tetapkanlah $\theta_n = \frac{\pi}{2n} + \pi\varepsilon_n$ dan $a_n =
\sin\theta_n$. Untuk $n$ yang besar, $\theta_n \in \intoo{0}{\frac\pi2}$
dan
\[
\theta_n - \theta_{n+1} = \frac{\pi}{2n(n+1)} +
\pi(\varepsilon_n - \varepsilon_{n+1})
= \frac{\pi}{2n^2} + O\Bigl(\frac{1}{n^3}\Bigr) > 0
\]
akhirnya, sehingga $(\theta_n)$ menurun ke $0$; dan karena $\sin$
naik pada $\intcc{0}{\frac\pi2}$, maka $(a_n)$ menurun ke $0$
pula. Jadi uji berselang-selingnya berlaku: konvergen --- tetapi tak mutlak,
karena $a_n \sim \frac{\pi}{2n}$.
\end{solution}

\begin{solution}{exo:b1:series:6}
Fungsi $f(t) = \frac{1}{t(\ln t)^\alpha}$ bersifat positif, \omterm{def:b1:continuity:continuous}{kontinu}, dan
menurun pada $\intco{2}{+\infty}$. Lalu dengan menyulihkan $u = \ln t$:
\[
\int_2^x \frac{\dd t}{t(\ln t)^\alpha}
= \int_{\ln 2}^{\ln x} \frac{\dd u}{u^\alpha},
\]
yang terbatas ketika $x \to \infty$ jika dan hanya jika $\alpha > 1$
(menurut perhitungan \cref{thm:b1:series:riemann}). Jadi menurut perbandingan
\omterm{thm:b1:integration:def}{integralnya}: konvergen jika dan hanya jika $\alpha > 1$. (Adapun \omterm{def:b1:series:def}{deret} \emph{bertipe Bertrand}
ini menunjukkan betapa halus \omterm{def:b1:topology:closure}{perbatasan} kekonvergenannya: karena $n\ln n$
divergen, sedangkan $n(\ln n)^{1.01}$ konvergen.)
\end{solution}

\begin{solution}{exo:b1:series:7}
Menurut batas garis singgung pada \cref{exo:b1:derivative:3} yang ditulis ulang
lewat ekspansi: untuk $x = \frac1n \in \intoc{0}{1}$, Taylor--Lagrange
bagi $\ln(1+x)$ pada orde $1$ memberikan $\ln(1 + x) = x - \frac{x^2}{2(1 +
c)^2}$ untuk suatu $c \in \intoo{0}{x}$, sehingga
\[
0 \leq u_n = \frac1n - \ln\Bigl(1 + \frac1n\Bigr) \leq
\frac{1}{2n^2} .
\]
Lalu perbandingan dengan \omterm{thm:b1:series:riemann}{deret Riemannnya}: $\sum u_n$ konvergen. Adapun jumlah
parsialnya berteleskop pada logaritmanya:
\[
\sum_{n=1}^{N} u_n = H_N - \ln(N+1)
\]
(karena $\sum_{n\leq N} \ln\frac{n+1}{n} = \ln(N+1)$). Jadi $H_N -
\ln(N+1)$ konvergen; lalu menambahkan $\ln\frac{N+1}{N} \to 0$, barisan
$H_N - \ln N$ konvergen. Adapun limitnya $\gamma \approx 0.5772$.
\end{solution}

\begin{solution}{exo:b1:series:8}
$\dfrac{1}{n(n+2)} = \dfrac{1/2}{n} - \dfrac{1/2}{n+2}$: jadi jumlah
parsialnya berteleskop dengan jeda $2$,
\[
S_N = \frac12\Bigl(1 + \frac12 - \frac{1}{N+1} -
\frac{1}{N+2}\Bigr) \longrightarrow \frac34 .
\]

Untuk $\sum \frac{n}{2^n}$: bagi $\abs x < 1$, dengan menurunkan jumlah
geometri yang hingga lalu beralih ke limitnya (karena semua \omterm{def:b1:series:def}{deret} di sini konvergen
mutlak, menurut \omterm{thm:b1:series:ratio}{uji nisbahnya}): dari $\sum_{n\geq0} x^n = \frac{1}{1-x}$,
kita memperoleh lewat perhitungan langsung dengan jumlah parsialnya
\[
\sum_{n=1}^{N} n x^{n-1}
= \frac{1 - (N+1)x^N + N x^{N+1}}{(1 - x)^2}
\xrightarrow[N\to\infty]{} \frac{1}{(1-x)^2}
\quad (\abs x < 1),
\]
(karena suku \omterm{def:b1:topology:closure}{perbatasannya} $N x^N \to 0$). Lalu di $x = \frac12$:
$\sum_{n\geq1} n\bigl(\frac12\bigr)^{n-1} = 4$, sehingga $\sum_{n\geq0}
\frac{n}{2^n} = \frac12 \times 4 = 2$.
\end{solution}

\begin{solution}{exo:b1:series:9}
Kelompokkanlah suku $\sum u_n$ dalam paket di antara pangkat $2$.
Adapun paket atasnya: untuk $2^k \leq n < 2^{k+1}$ ada $2^k$ suku,
yang masing-masingnya $\leq u_{2^k}$:
\[
\sum_{n=1}^{2^{K+1}-1} u_n
= \sum_{k=0}^{K} \sum_{n=2^k}^{2^{k+1}-1} u_n
\leq \sum_{k=0}^{K} 2^k u_{2^k} .
\]
Adapun paket bawahnya: setiap suku paket yang sama bernilai $\geq u_{2^{k+1}}$,
sehingga $\sum_{n=2^k}^{2^{k+1}-1} u_n \geq 2^k u_{2^{k+1}} = \frac12
\cdot 2^{k+1} u_{2^{k+1}}$, jadi
\[
\sum_{n=1}^{2^{K+1}-1} u_n \geq \frac12 \sum_{k=1}^{K+1} 2^{k}
u_{2^{k}} .
\]
Kedua perbandingan jumlah parsialnya berjalan ke kedua arahnya (karena sukunya taknegatif,
\cref{thm:b1:series:positive}): sehingga kedua \omterm{def:b1:series:def}{deretnya} bersifat sama.

Untuk Riemann: $u_n = n^{-\alpha}$ memberikan $2^k u_{2^k} =
2^{k(1-\alpha)}$, yaitu \omterm{ex:b1:series:geometric}{deret geometri}, yang konvergen jika dan hanya jika $2^{1 - \alpha}
< 1$ jika dan hanya jika $\alpha > 1$. Adapun untuk Bertrand
(\cref{exo:b1:series:6}): $u_n = \frac{1}{n(\ln n)^\alpha}$ memberikan
$2^k u_{2^k} = \frac{1}{(k\ln 2)^\alpha}$, yaitu \omterm{thm:b1:series:riemann}{deret Riemann} dalam $k$:
jadi konvergen jika dan hanya jika $\alpha > 1$.
\end{solution}

\begin{solution}{exo:b1:series:10}
Lewat kesamaan geometri yang hingga dengan nisbah $-t^2$:
\[
\frac{1}{1 + t^2} = \sum_{k=0}^{n-1} (-1)^k t^{2k}
+ \frac{(-1)^n t^{2n}}{1 + t^2} .
\]
Lalu integralkanlah atas $\intcc{0}{1}$ (karena ruas kirinya terintegralkan menjadi $\arctan
1 = \frac\pi4$, \cref{prop:b1:functions:arcderiv}):
\[
\frac{\pi}{4} = \sum_{k=0}^{n-1} \frac{(-1)^k}{2k+1}
+ (-1)^n \int_0^1 \frac{t^{2n}}{1+t^2}\,\dd t,
\qquad
0 \leq \int_0^1 \frac{t^{2n}}{1+t^2}\,\dd t \leq \int_0^1 t^{2n}\dd
t = \frac{1}{2n+1} .
\]
Lalu membiarkan $n \to \infty$ membuktikan rumusnya, dan batas yang terpampang
atas \omterm{thm:b1:integration:def}{integralnya} persis merupakan batas sisanya: sehingga setelah menjumlahkan sampai
$N$ (yakni $n = N + 1$ suku), $\abs{R_N} \leq \frac{1}{2N + 3}$.
\end{solution}

\begin{solution}{exo:b1:series:11}
$n^{1/n} = \eu^{\frac{\ln n}{n}} \to \eu^0 = 1$
(\cref{prop:b1:functions:powerrules}). Sehingga
\[
\frac{1}{n^{1 + 1/n}} = \frac{1}{n}\,\eu^{-\frac{\ln n}{n}}
\sim \frac{1}{n} ,
\]
lalu uji kesetaraannya (\cref{thm:b1:series:positive})
membandingkannya dengan \omterm{def:b1:series:def}{deret} harmonik yang divergen: jadi \emph{divergen},
meskipun setiap eksponennya $1 + \frac1n$ melampaui $1$. Adapun kriteria
Riemann menyangkut eksponen $\alpha$ yang \emph{tetap}; sedangkan eksponen
yang meluncur turun ke $1$ dapat kehilangan seluruh kelonggarannya, seperti di sini.
\end{solution}

\begin{solution}{exo:b1:series:12}
Misalkan $\varepsilon > 0$. Menurut kriteria Cauchy bagi
\omterm{def:b1:series:def}{deret} yang konvergen (\cref{thm:b1:seq:complete} yang diterapkan pada
jumlah parsialnya), ada $N$ dengan $\sum_{k=n+1}^{2n} u_k \leq
\varepsilon$ untuk $n \geq N$. Lalu menurut kemonotonannya masing-masing dari $n$
sukunya bernilai $\geq u_{2n}$:
\[
n\,u_{2n} \leq \sum_{k=n+1}^{2n} u_k \leq \varepsilon
\quad\Longrightarrow\quad 2n\,u_{2n} \leq 2\varepsilon ,
\]
dan untuk indeks yang ganjil $(2n+1)\,u_{2n+1} \leq (2n+1)\,u_{2n} \leq
2\bigl(2n\,u_{2n}\bigr) \leq 4\varepsilon$ untuk $n \geq N$: jadi pada
kedua keparitasannya, $n u_n \to 0$.

Adapun konversnya salah: karena $u_n = \frac{1}{n\ln n}$ mempunyai $n u_n =
\frac{1}{\ln n} \to 0$, namun \omterm{def:b1:series:def}{deretnya} divergen
(\cref{exo:b1:series:6}, dengan $\alpha = 1$). Dan kemonotonannya diperlukan:
misalkan $u_n = \frac1n$ ketika $n$ kuadrat sempurna dan $u_n =
2^{-n}$ selainnya: maka \omterm{def:b1:series:def}{deretnya} konvergen (karena suku kuadratnya berjumlah
seperti $\sum \frac{1}{k^2}$, dan sisanya secara geometri), tetapi $n u_n =
1$ sepanjang kuadratnya.
\end{solution}

\begin{solution}{pb:b1:series:1}
\textbf{1.} Di sini $a_{n+1} - a_n = \frac{1}{n+1} - \ln\frac{n+1}{n}
\leq 0$ karena $\ln(1 + \frac1n) \geq \frac{1/n}{1 + 1/n} =
\frac{1}{n+1}$; dan $b_{n+1} - b_n = \frac{1}{n+1} -
\ln\frac{n+2}{n+1} \geq 0$ karena $\ln(1 + \frac{1}{n+1}) \leq
\frac{1}{n+1}$. Adapun celahnya $a_n - b_n = \ln(1 + \frac1n) \to 0$:
jadi berdampingan (\cref{thm:b1:seq:adjacent}), dengan limit bersamanya $\lim
a_n = \gamma$ (\cref{exo:b1:series:7}), dan $b_n \leq \gamma
\leq a_n$.

\textbf{2.} Di sini $b_{10} = 2.928968 - \ln 11 = 0.5311$ dan $a_{10} =
2.928968 - \ln 10 = 0.6264$: sehingga $\gamma \in
\intcc{0.5311}{0.6264}$. Adapun celahnya $\ln 1.1 \approx 0.095$ dan
menyusut seperti $\frac1n$: sehingga empat desimal (dengan $\text{celah} \leq
10^{-4}$) akan memerlukan $n \approx 10^4$ --- jadi apitannya
benar tetapi lamban.

\textbf{3.} Lewat teleskop $a_n - a_{m+1} = \sum_{k=n}^{m} w_k$ lalu
membiarkan $m \to \infty$: $a_n - \gamma = \sum_{k\geq n} w_k$
(yaitu limit jumlah parsialnya). Lebih lanjut
\[
w_k = \int_k^{k+1} \frac{\dd t}{t} - \frac{1}{k+1}
= \int_k^{k+1} \Bigl(\frac1t - \frac{1}{k+1}\Bigr)\dd t
= \int_k^{k+1} \frac{k + 1 - t}{t\,(k+1)}\,\dd t .
\]
Pada $\intcc{k}{k+1}$: $\frac{1}{(k+1)^2} \leq \frac{1}{t(k+1)}
\leq \frac{1}{k(k+1)}$, dan $\int_k^{k+1}(k + 1 - t)\dd t =
\frac12$: sehingga $\frac{1}{2(k+1)^2} \leq w_k \leq
\frac{1}{2k(k+1)}$.

\textbf{4.} Untuk yang atas: $\sum_{k \geq n} \frac{1}{2k(k+1)} =
\frac12\sum_{k\geq n}\bigl(\frac1k - \frac{1}{k+1}\bigr) =
\frac{1}{2n}$ (secara teleskopis). Untuk yang bawah: $\frac{1}{2(k+1)^2} \geq
\frac{1}{2(k+1)(k+2)}$, yang jumlahnya berteleskop menjadi
$\frac{1}{2(n+1)}$. Digabung dengan pertanyaan 3:
\[
\frac{1}{2(n+1)} \leq H_n - \ln n - \gamma \leq \frac{1}{2n} .
\]

\textbf{5.} Kurangkanlah $\frac{1}{2n}$: maka $\gamma_n - \gamma \in
\intcc{\frac{1}{2(n+1)} - \frac{1}{2n}}{0} =
\intcc{-\frac{1}{2n(n+1)}}{0}$: jadi taksiran yang terkoreksi itu persis
sampai $O\bigl(\frac{1}{n^2}\bigr)$, dan selalu dari bawah.

\textbf{6.} Di sini $\gamma_{100} = 5.1873775 - \ln 100 - 0.005 =
0.5772073$, dengan $0 \leq \gamma - \gamma_{100} \leq
\frac{1}{20200} < 5\cdot10^{-5}$: sehingga $0.577207 \leq \gamma
\leq 0.577257$, yakni $\gamma = 0.5772 \pm 5\cdot10^{-5}$ (dengan nilai
sejatinya $0.5772156\dots$) --- jadi empat desimal bersertifikat dari
seratus suku, terhadap sepuluh ribu bagi pertanyaan 2.

\textbf{7.} Adapun dividen pertamanya:
\[
H_{2m} - H_m = \Bigl(\ln 2m + \gamma + \frac{1}{4m}\Bigr)
- \Bigl(\ln m + \gamma + \frac{1}{2m}\Bigr) +
O\Bigl(\frac{1}{m^2}\Bigr)
= \ln 2 - \frac{1}{4m} + O\Bigl(\frac{1}{m^2}\Bigr).
\]
Untuk yang kedua: kebalikan bilangan genap sampai $2m$ berjumlah $\frac12 H_m$, sehingga
$\sum_{j=1}^{m}\frac{1}{2j-1} = H_{2m} - \frac12 H_m = \frac12
\ln m + \ln 2 + \frac\gamma2 + o(1)$, dan $\sum_{j=1}^m
\frac{1}{2j} = \frac12\ln m + \frac\gamma2 + o(1)$.

\textbf{8.} Dengan memisahkan suku genapnya dua kali:
$\sum_{k=1}^{2m}\frac{(-1)^{k-1}}{k} = H_{2m} - 2\cdot\frac12
H_m = H_{2m} - H_m$. Lalu menurut pertanyaan 7 ini sama dengan $\ln 2 -
\frac{1}{4m} + O(m^{-2})$: sehingga jumlahnya $\ln 2$
(\cref{ex:b1:series:ln2} lagi) \emph{beserta} kecepatannya.

\textbf{9.} Untuk $N = 2m$: $S - S_N = \frac{1}{4m} + O(m^{-2}) =
\frac{1}{2N} + O(N^{-2})$. Sedangkan untuk $N = 2m + 1$: $S_{N} = S_{2m} +
\frac{1}{2m+1}$, sehingga
\[
S - S_N = \Bigl(\frac{1}{4m} - \frac{1}{2m+1}\Bigr) +
O\Bigl(\frac{1}{m^2}\Bigr)
= -\frac{1}{4m} + O\Bigl(\frac{1}{m^2}\Bigr)
= -\frac{1}{2N} + O\Bigl(\frac{1}{N^2}\Bigr).
\]
Jadi pada kedua kasusnya $S - S_N \sim \frac{(-1)^N}{2N}$: yakni separuh batas
kasus terburuknya $a_{N+1}$, dengan tanda berselang-seling yang diketahui.

\textbf{10.} Perataannya membunuh suku utamanya yang berayun:
\[
S - \tilde S_N = \frac{(S - S_N) + (S - S_{N+1})}{2}
= \frac{(-1)^N}{2}\Bigl(\frac{1}{2N} - \frac{1}{2(N+1)}\Bigr)
+ O\Bigl(\frac{1}{N^2}\Bigr) = O\Bigl(\frac{1}{N^2}\Bigr).
\]
Secara numerik: $S_{10} = 0.645635$, $S_{11} = 0.736544$, $\tilde
S_{10} = 0.691089$, dan $\ln 2 = 0.693147$: sehingga galatnya turun dari
$4.8\cdot10^{-2}$ menjadi $2.1\cdot10^{-3}$ --- jadi satu penjumlahan, dua puluh
kali lebih baik.

\textbf{11.} Galat berselang-selingnya berubah tanda pada setiap langkahnya,
sehingga jumlah parsial yang berurutan mengangkangi limitnya dan rata-ratanya
meniadakan suku orde pertamanya; sedangkan galat apitannya $H_n - \ln n -
\gamma \approx \frac{1}{2n}$ bertanda tetap, sehingga tak ada perataan
sepanjang $n$ yang dapat meniadakannya. Adapun merata-ratakan kedua \emph{apitannya} memang
membantu: $\frac{a_n + b_n}{2} = H_n - \ln\sqrt{n(n+1)}$, dan karena
$\ln\sqrt{n(n+1)} = \ln\bigl(n + \frac12\bigr) + O(n^{-2})$,
\[
H_n - \ln\Bigl(n + \frac12\Bigr)
= \Bigl(H_n - \ln n - \frac{1}{2n}\Bigr) + \frac{1}{8n^2} +
O\Bigl(\frac{1}{n^3}\Bigr) = \gamma + O\Bigl(\frac{1}{n^2}\Bigr)
\]
(menurut pertanyaan 5 dan $\ln(1 + \frac{1}{2n}) = \frac{1}{2n} -
\frac{1}{8n^2} + O(n^{-3})$). Periksalah di $n = 10$: $H_{10} -
\ln 10.5 = 0.57759$, yang sudah berada dalam jarak $4\cdot10^{-4}$ dari
$\gamma$.

\textbf{12.} Di sini $\sum_{j\leq m} \frac{1}{2j-1} \geq \sum_{j \leq m}
\frac{1}{2j} = \frac12 H_m \to +\infty$: sehingga baik bagian positif maupun
bagian negatif \omterm{def:b1:series:def}{deret} harmonik berselang-selingnya divergen.

\textbf{13.} Tulislah $u_n^\pm = \frac{\abs{u_n} \pm u_n}{2} \geq
0$, sehingga $u_n = u_n^+ - u_n^-$ dan $\abs{u_n} = u_n^+ + u_n^-$. Seandainya
$\sum u_n^+$ konvergen, maka $\sum u_n^- = \sum (u_n^+ - u_n)$
akan konvergen (sebagai selisih \omterm{def:b1:series:def}{deret} yang konvergen), sehingga $\sum
\abs{u_n}$ pun demikian: yang bertentangan dengan kekonvergenan bersyaratnya. Jadi menurut
kesimetriannya kedua $\sum u_n^\pm$ divergen (ke $+\infty$): yakni sebuah
cadangan massa positif dan negatif yang tak hingga.

\textbf{14.} Misalkan $S = \sum u_n$, $\varepsilon > 0$, dan $N$ dengan
$\sum_{n > N} \abs{u_n} \leq \varepsilon$ (menurut kriteria Cauchy bagi
$\sum\abs{u_n}$). Misalkan $M_0$ cukup besar sehingga $\sigma(\{0,
\dots, M_0\}) \supseteq \{0, \dots, N\}$. Untuk $M \geq M_0$, selisih
$\sum_{m \leq M} u_{\sigma(m)} - \sum_{n \leq N} u_n$
merupakan jumlah hingga atas suku $u_n$ yang \emph{berbeda} dengan $n > N$,
sehingga bernilai mutlak $\leq \varepsilon$; dan $\abs{S -
\sum_{n\leq N} u_n} \leq \varepsilon$ pula. Jadi jumlah parsial yang tersusun
ulang itu akhirnya berada dalam jarak $2\varepsilon$ dari $S$:
sehingga $\sum u_{\sigma(n)} = S$. Jadi \omterm{thm:b1:series:absolute}{kekonvergenan mutlak} bersifat
kebal terhadap penyusunan ulang.

\textbf{15.} Setiap fase prosedur rakusnya berakhir setelah
berhingga banyak suku, karena suku positif (masing-masing negatif)
yang tersisa saja sudah berjumlah parsial divergen
(pertanyaan 12): sehingga jumlah berjalannya akhirnya pasti melintasi $t$. Jadi
prosedurnya berselang-seling melewati tak hingga banyak fase yang hingga,
sambil memakai suku positifnya berurutan dan suku negatifnya
berurutan: sehingga setiap sukunya terpakai tepat sekali --- yakni sebuah penyusunan ulang.
Setelah persilangan pertamanya, di antara dua persilangan yang berurutan
jumlah parsialnya bergerak secara monoton menuju $t$, dan pada sebuah persilangan
ia melampaui paling banyak sebesar suku yang baru ditambahkan; lalu karena suku
yang terpakai pada persilangan ke-$j$ berindeks sekurang-kurangnya $j$ pada
kelasnya, maka pelampauan itu menuju $0$. Jadi jumlah parsialnya
konvergen ke $t$: sehingga setiap bilangan real merupakan jumlah suatu
penyusunan ulang.

\textbf{16.} Slot positifnya menerima $\frac{1}{2j-1}$ untuk
$j = 1, 2, \dots$ secara berurutan, dan slot negatifnya $\frac{1}{2j}$ secara
berurutan: sehingga setiap suku \omterm{def:b1:series:def}{deret} harmonik berselang-selingnya muncul
tepat sekali. Adapun untuk $(p, q) = (1, 2)$, bloknya adalah
$\bigl(1, -\frac12, -\frac14\bigr)$, $\bigl(\frac13, -\frac16,
-\frac18\bigr)$, $\bigl(\frac15, -\frac1{10},
-\frac1{12}\bigr)$, \dots\ --- yaitu \omterm{def:b1:series:def}{deret} yang terpampang.

\textbf{17.} Karena $\frac{1}{4k-2} = \frac12\cdot\frac{1}{2k-1}$:
\[
\frac{1}{2k-1} - \frac{1}{4k-2} - \frac{1}{4k}
= \frac12\,\frac{1}{2k-1} - \frac12\,\frac{1}{2k}
= \frac12\Bigl(\frac{1}{2k-1} - \frac{1}{2k}\Bigr).
\]
Lalu menjumlahkannya atas $k = 1, \dots, K$: $T_{3K} = \frac12
\sum_{k=1}^{K}\bigl(\frac{1}{2k-1} - \frac{1}{2k}\bigr) =
\frac12 S_{2K}$: jadi pada setiap jumlah parsial ketiganya, \omterm{def:b1:series:def}{deret}
yang tersusun ulang itu \emph{persis} separuh yang aslinya.

\textbf{18.} Setelah $K$ blok yang lengkap, jumlah parsial yang
tersusun ulang adalah $\sum_{j=1}^{pK}\frac{1}{2j-1} -
\sum_{j=1}^{qK}\frac{1}{2j}$, lalu pertanyaan 7 menilainya:
\[
\Bigl(\frac{\ln(pK)}{2} + \ln 2 + \frac\gamma2\Bigr)
- \Bigl(\frac{\ln(qK)}{2} + \frac\gamma2\Bigr) + o(1)
= \ln 2 + \frac12\ln\frac pq + o(1) :
\]
sehingga $\gamma$-nya saling meniadakan, $\ln K$-nya saling meniadakan, dan nisbah
$\frac pq$-nya bertahan.

\textbf{19.} Sebuah jumlah parsial di dalam blok $K + 1$ berselisih dari
jumlah $K$ bloknya paling banyak $p + q$ suku, yang masing-masingnya bernilai mutlak
kira-kira $\leq \frac{1}{2qK}$, sehingga berselisih $O\bigl(\frac1K\bigr) \to 0$:
jadi seluruh barisan jumlah parsialnya berlimit sama $\ln 2 +
\frac12\ln\frac pq$. Adapun untuk $(1, 2)$: $\ln 2 + \frac12\ln\frac12 =
\frac{\ln 2}{2} = 0.34657\dots$, dan memang $T_9 = 0.30833$
merayap menujunya: karena menurut pertanyaan 17, $T_{3K} = \frac12 S_{2K}$
konvergen dengan persis separuh galat harmonik berselang-selingnya.
Jadi suku yang sama, separuh jumlahnya.

\textbf{20.} Untuk $(1,1)$: $\ln 2 + \frac12\ln 1 = \ln 2$ --- yaitu
urutan aslinya, jadi taat asas. Untuk $(2,1)$: $\frac32\ln 2 \approx
1.0397$. Adapun menu $(p,q)$-nya tepat mencapai keluarga
terbilang yang \omterm{def:b1:topology:dense}{padat} $\ln 2 + \frac12\ln r$, dengan $r \in \Q_{>0}$; sedangkan resep
rakus Riemann (pertanyaan 15) mencapai \emph{setiap} bilangan real. Jadi struktur membeli
rumus; sedangkan kerakusan membeli ketotalan.

\textbf{21.} Jumlah parsialnya berteleskop: $\sum_{k=1}^{N}
\bigl(\frac1k - \ln\frac{k+1}{k}\bigr) = H_N - \ln(N+1) = b_N
\to \gamma$: sehingga \omterm{def:b1:series:def}{deret} pada \cref{exo:b1:series:7} berjumlah persis
\omterm{pb:b1:series:1}{konstanta Euler}.

\textbf{22.} Rakus bagi $t = 1$: suku positif pertamanya membawa
jumlahnya persis ke $1$, bukan melampauinya, sehingga sebuah positif kedua
diambil untuk melintasinya: $1, \frac13$ (dengan jumlah $1.3333 > 1$), lalu
$-\frac12$ ($0.8333$), $\frac15$ ($1.0333$), $-\frac14$
($0.7833$), $\frac17, \frac19$ ($1.0373$), $-\frac16$
($0.8706$), $\frac1{11}, \frac1{13}$ ($1.0384$), $-\frac18$
($0.9134$), $\frac1{15}$ ($0.9801$), \dots\ --- jadi jumlahnya bernapas
di sekitar $1$ dengan amplitudo yang kian kecil, dengan dua positif kini
diperlukan per daur karena negatifnya lebih besar.

\textbf{23.} Bangunlah bloknya: pada tahap $j$, tambahkanlah cukup banyak suku
positif yang belum terpakai untuk menaikkan jumlah parsialnya sekurang-kurangnya $1$
(yang mungkin: karena suku positif yang tersisa berjumlah divergen,
pertanyaan 12), lalu tambahkanlah satu suku negatif tunggal $-\frac{1}
{2j}$. Maka setiap suku positifnya akhirnya terpakai (karena setiap tahapnya memakai
sekurang-kurangnya satu), dan setiap yang negatif pun demikian (satu per tahap): jadi sebuah
penyusunan ulang. Adapun setiap tahapnya mengubah jumlahnya sebesar $\geq 1 -
\frac{1}{2j} \geq \frac12$: sehingga jumlah parsialnya melampaui
$\frac{j}{2}$ setelah tahap $j$ dan pertambahan di dalam sebuah tahapnya
positif kecuali yang terakhir, yang terbatas oleh $\frac{1}{2j} \to 0$:
jadi divergen ke $+\infty$.

\textbf{24.} Menurut hukumnya, $H_N \geq 20 \iff \ln N \geq 20 -
\gamma - \frac{1}{2N} + O(N^{-2})$: sehingga ambangnya $N^*$
memenuhi $\ln N^* = 20 - \gamma + o(1)$, yakni $N^* =
\eu^{20 - \gamma}(1 + o(1)) \approx \eu^{19.4228} \approx
2.72\cdot10^{8}$ --- yang di dalam jendela kasar
$\intcc{1.8\cdot10^8}{4.9\cdot10^8}$ pada
\cref{ex:b1:series:harmonicstack}, dan dipakukan oleh $\gamma$.

\textbf{25.} (i) Pada $H_n = \ln n + \gamma + \frac{1}{2n} +
O(n^{-2})$: $\ln n$-nya adalah \omterm{thm:b1:integration:def}{integralnya}, $\gamma$-nya harga
mengganti sebuah jumlah dengan sebuah \omterm{thm:b1:integration:def}{integral} (yaitu konstanta analisis yang
sungguh baru), dan $\frac{1}{2n}$-nya koreksi pertamanya --- yakni
bayangan trapesiumnya. (ii) Kekonvergenan bersyarat bersandar pada
peniadaan di antara dua cadangan yang tak hingga (pertanyaan 13), sehingga
penyusunan ulangnya menimbang ulang cadangannya; sedangkan \omterm{thm:b1:series:absolute}{kekonvergenan mutlak} mempunyai
massa total yang hingga, dan taksiran ekor pertanyaan 14 bersifat
buta urutan. (iii) Adapun jumlah $(p,q)$-nya \emph{dihitung}: karena
hukum $\gamma$-nya mengubah setiap jumlah parsial yang tersusun ulang menjadi $\ln 2 +
\frac12\ln\frac pq + o(1)$, dengan $\gamma$-nya sendiri saling meniadakan ---
jadi latihan pembukuan asimtotik, bukan argumen yang abstrak.
(iv) Berikutnya: hasil kali dan keterjumlahan takbersyarat bagi \omterm{def:b1:series:def}{deret}
pangkat pada jilid Tahun ke-2; dan soal akhir pekan jilid Tahun ke-3
tentang rumus Stirling, yang di situ pembukuan jumlah lawan \omterm{thm:b1:integration:def}{integralnya},
bila didorong satu orde lebih jauh, menghasilkan $\sqrt{2\pi}$
sendiri.

\end{solution}
